\documentclass[10pt]{article}
\providecommand{\FIGDIR}{fig}
\newcommand{\DIGESTDATE}{21 August 2026}
% ===================== PM Research Watch - shared preamble =====================
\usepackage[T1]{fontenc}
\usepackage[utf8]{inputenc}
\usepackage{amsmath,amssymb}
\usepackage{tgpagella}
\usepackage{tgheros}
\renewcommand{\sfdefault}{qhv}
\renewcommand{\ttdefault}{lmtt}
\usepackage{microtype}
\usepackage[english]{babel}

\usepackage[a4paper,top=17mm,bottom=17mm,left=17mm,right=17mm,
            headsep=6mm,footskip=10mm]{geometry}
\usepackage{graphicx}
\usepackage{xcolor}
\usepackage{colortbl}
\usepackage{tikz}
\usetikzlibrary{calc}
\usepackage[most]{tcolorbox}
\usepackage{titlesec}
\usepackage{fancyhdr}
\usepackage{booktabs}
\usepackage{array}
\usepackage{longtable}
\usepackage{enumitem}
\usepackage{multicol}
\usepackage{ragged2e}
\usepackage[hidelinks]{hyperref}

% ------------------------------------------------------------------ palette
\definecolor{Ink}    {HTML}{14262E}
\definecolor{Deep}   {HTML}{0F4C5C}
\definecolor{Teal}   {HTML}{1F7A8C}
\definecolor{Sky}    {HTML}{4FA3C4}
\definecolor{Sage}   {HTML}{6FA287}
\definecolor{Amber}  {HTML}{E3A008}
\definecolor{Clay}   {HTML}{B5643C}
\definecolor{Coral}  {HTML}{D1495B}
\definecolor{Violet} {HTML}{6A5B9E}
\definecolor{Slate}  {HTML}{4A6572}
\definecolor{Paper}  {HTML}{FBF9F5}
\definecolor{Mist}   {HTML}{EEEAE2}
\definecolor{Rule}   {HTML}{D9D3C9}

\pagecolor{Paper}
\color{Ink}

\hypersetup{colorlinks=true, linkcolor=Deep, urlcolor=Teal, citecolor=Teal,
            pdftitle={PM Research Watch}, pdfauthor={Automated daily digest}}

% ------------------------------------------------------------------ headings
\titleformat{\section}
  {\sffamily\bfseries\large\color{Deep}}{}{0pt}
  {}[{\vspace{1.2mm}\color{Amber}\titlerule[1.1pt]}]
\titlespacing*{\section}{0pt}{5.5mm}{2.8mm}

\titleformat{\subsection}
  {\sffamily\bfseries\normalsize\color{Teal}}{}{0pt}{}
\titlespacing*{\subsection}{0pt}{3.6mm}{1.4mm}

\setlist[itemize]{leftmargin=4.2mm, itemsep=1.1mm, topsep=1.3mm, parsep=0pt,
                  label=\textcolor{Teal}{\footnotesize$\blacktriangleright$}}

% ------------------------------------------------------------------ header/footer
\pagestyle{fancy}
\fancyhf{}
\renewcommand{\headrulewidth}{0pt}
\renewcommand{\footrulewidth}{0pt}
\fancyhead[L]{\sffamily\scriptsize\color{Slate}PM RESEARCH WATCH\;\textcolor{Rule}{\textbar}\;\DIGESTDATE}
\fancyhead[R]{\sffamily\scriptsize\color{Slate}Particulate matter monitoring \& health effects}
\fancyfoot[C]{\sffamily\scriptsize\color{Slate}\thepage}

% ------------------------------------------------------------------ components
% metric tile
\newcommand{\metric}[3]{%
  \begin{tcolorbox}[enhanced, colback=white, colframe=Rule, boxrule=0.5pt,
    arc=1.6mm, left=2.2mm, right=2.2mm, top=1.8mm, bottom=1.8mm,
    borderline west={1.7mm}{0pt}{#3}, nobeforeafter, valign=center,
    height=17mm, width=\linewidth]
    \RaggedRight
    {\sffamily\bfseries\fontsize{17}{17}\selectfont\color{#3}#1}\\[0.5mm]
    {\sffamily\fontsize{6.5}{7.4}\selectfont\color{Slate}#2}
  \end{tcolorbox}}

% pull-quote / key insight box
\newtcolorbox{keybox}[1][]{enhanced, colback=Mist, colframe=Mist,
  arc=1.6mm, left=3.2mm, right=3.2mm, top=2.4mm, bottom=2.4mm,
  borderline west={1.7mm}{0pt}{Amber}, #1}

% subtopic band
\newcommand{\band}[2]{%
  \vspace{3.2mm}\noindent
  \begin{tcolorbox}[enhanced, colback=#2, colframe=#2, arc=1.2mm,
    left=2.6mm, right=2.6mm, top=1.1mm, bottom=1.1mm, boxrule=0pt,
    width=\linewidth, nobeforeafter]
    {\sffamily\bfseries\small\color{white}#1}
  \end{tcolorbox}\vspace{1.6mm}}

% one paper entry: {short cite}{title}{journal}{doi}{body}
\newcommand{\paperentry}[5]{%
  \par\vspace{1.9mm}\noindent
  {\sffamily\bfseries\footnotesize\color{Ink}#1}\;%
  {\sffamily\fontsize{7}{8}\selectfont\color{Slate}#3}\\[0.4mm]
  {\footnotesize\itshape\color{Deep}#2}\\[0.6mm]
  {\footnotesize #4}\\[0.3mm]
  {\fontsize{6.6}{7.6}\selectfont\color{Slate}\texttt{doi:}\,\href{https://doi.org/#5}{\texttt{\detokenize{#5}}}}
  \par}

% caption style
\newcommand{\figcap}[1]{\par\vspace{1.0mm}
  {\fontsize{7.2}{8.6}\selectfont\color{Slate}\textsf{\textbf{Fig.}}\; #1\par}\vspace{1.5mm}}

\providecommand{\permil}{\textperthousand}

\setlength{\parindent}{0pt}
\setlength{\parskip}{0pt}
\renewcommand{\arraystretch}{1.18}


\begin{document}

% ==================================================================== MASTHEAD
\thispagestyle{empty}
\begin{tikzpicture}[remember picture, overlay]
  \fill[Deep] (current page.north west) rectangle
      ($(current page.north east)+(0,-67mm)$);
  \fill[Amber] ($(current page.north west)+(0,-67mm)$) rectangle
      ($(current page.north east)+(0,-68.4mm)$);
  \foreach \x/\y/\r in {22/12/0.9, 41/26/1.5, 63/9/0.7, 88/31/1.1, 112/17/1.9,
                        138/29/0.8, 159/13/1.3, 178/54/1.0, 196/21/1.6,
                        30/44/0.6, 74/57/1.2, 125/48/0.9, 168/6/0.8, 191/59/0.7,
                        52/18/1.0, 100/60/0.6, 148/20/1.1, 66/49/1.4, 205/45/0.8}
    \fill[Sky, opacity=0.30] ($(current page.north west)+(\x mm,-\y mm)$) circle (\r mm);
\end{tikzpicture}

\vspace*{4mm}
{\sffamily\fontsize{7.6}{9}\selectfont\color{Sky}%
 AUTOMATED DAILY LITERATURE DIGEST\;\textbullet\;PhD RESEARCH WATCH\par}
\vspace{2.2mm}
{\sffamily\bfseries\fontsize{27}{29}\selectfont\color{white}%
 Particulate Matter\par}
{\sffamily\fontsize{27}{29}\selectfont\color{Sky}%
 Monitoring \& Health Effects\par}
\vspace{5.0mm}
{\sffamily\fontsize{8.4}{10}\selectfont\color{white}%
 \mbox{Issue 2026-08-21}\;\;\textcolor{Amber}{\textbar}\;\;\mbox{Entry window 21 Aug 2026}\;\;
 \textcolor{Amber}{\textbar}\;\;\mbox{19 in scope, 46 screened out}\par}
\vspace{18mm}

% -------------------------------------------------------------- metric strip
\noindent
\begin{minipage}[t]{0.190\linewidth}\metric{19}{RECORDS IN SCOPE\\(46 SCREENED OUT)}{Deep}\end{minipage}\hfill
\begin{minipage}[t]{0.190\linewidth}\metric{7}{OF 10 CLUSTERS\\OCCUPIED}{Teal}\end{minipage}\hfill
\begin{minipage}[t]{0.190\linewidth}\metric{8}{EXPOSURE /\\MODELLING (42\%)}{Sky}\end{minipage}\hfill
\begin{minipage}[t]{0.190\linewidth}\metric{3}{PREPRINTS\\THIS ISSUE}{Amber}\end{minipage}\hfill
\begin{minipage}[t]{0.190\linewidth}\metric{1}{RISK ESTIMATE\\ON THE PLOT}{Clay}\end{minipage}

\vspace{6mm}

% -------------------------------------------------------------- signal of day
\section*{\sffamily\bfseries\color{Deep}Signal of the day}
\vspace{-1.5mm}{\color{Amber}\rule{\linewidth}{1.1pt}}\vspace{2.5mm}

\begin{keybox}
{\sffamily\bfseries\small\color{Deep}Three independent papers today answer the same
question --- \emph{which part of the mass?} --- and none of them answers it with the mass.}\\[1.4mm]
\footnotesize
Wei et al.\ (\textit{Environ Res}) ran generalized weighted quantile sum regression over
\textbf{15 PM$_{2.5}$ chemical components} in 528{,}127 SEER-Medicare patients aged
$\geq$65 with lung cancer, and report \textbf{RR 1.011 (95\% CI 1.010--1.013) per decile
increase in the mixture}, with silicon, nitrate, vanadium, zinc and iron carrying the
largest weights --- a road-dust, traffic and heavy-fuel-oil signature rather than a mass
signature. Chen et al.\ (\textit{Environ Toxicol Pharmacol}) decompose PM$_{2.5}$ into
five constituents in 24{,}940 UK Biobank participants and find \emph{ammonium} the
strongest single predictor of incident depression, with 140 plasma proteins and 3
metabolites mediating up to 8.0\%. And Chai et al.\ (\textit{Environ Pollut}) show why
that matters upstream: in suburban Nanjing, secondary inorganic aerosol is
\textbf{49.55\% of PM$_{2.5}$ mass}, nitrate-dominated and ammonia-buffered, so
\emph{PM$_{2.5}$ mass barely responds to moderate NH$_3$ cuts} --- it takes $>$60\%
reduction to move, whereas cutting HNO$_3$ alone performs as well as cutting both.
The common implication for instrumentation is direct: a monitor that reports only mass
cannot distinguish the exposure that carries the risk from the exposure that does not,
and cannot tell a policy-maker which precursor to regulate.
\end{keybox}

\vspace{3mm}

% -------------------------------------------------------------- exec brief
\section*{\sffamily\bfseries\color{Deep}What changed today}
\vspace{-1.5mm}{\color{Amber}\rule{\linewidth}{1.1pt}}\vspace{2.5mm}

\footnotesize
\begin{itemize}
\item \textbf{The issue is unusually weighted toward exposure and deposition.} Eight of
nineteen records (42\%) sit in \textit{Exposure assessment \& modelling}, and four of
those eight measure PM by what it \emph{lands on} rather than what it does in air ---
pine needles in Milazzo, privet leaves after two Albanian fires, ten tree species in
Delhi, and Elsevier's aerosol--precipitation pair. Deposition biomonitoring is cheap,
retrospective and spatially dense; it is also unstandardised, and no two of these four
papers report leaf-loading in the same units.

\item \textbf{An airborne-microplastic baseline for Southeast Asia.} Kan et al.\ report
the first Phnom Penh measurements: \textbf{0.97 $\pm$ 0.30 MP\,m$^{-3}$ in TSP and
0.33 $\pm$ 0.10 MP\,m$^{-3}$ in the PM$_{2.5}$ fraction}, 82\% below 30\,$\mu$m Feret
diameter. Roughly a third of the airborne microplastic count sits inside the size
fraction that conventional PM$_{2.5}$ monitoring already samples --- and that no
conventional PM$_{2.5}$ monitor identifies.

\item \textbf{An ammonia-rich regime where the obvious control lever does not work.}
The Nanjing thermodynamic sensitivity analysis is total-nitrate-limited, not
ammonia-limited. This is a concrete counter-example to the assumption that NH$_3$
abatement is a general-purpose PM$_{2.5}$ control.

\item \textbf{Indoor infiltration quantified at a useful sample size.} A preprint
covering \textbf{48 residences across eight northern Chinese cities}, four seasons,
32{,}256 hourly paired indoor/outdoor observations, gives an infiltration factor of
\textbf{0.28 $\pm$ 0.15} and attributes \textbf{57--69\% of indoor PM$_{2.5}$ to outdoor
sources}, peaking in the heating season. Air-purifier operation was the single strongest
determinant in the model ($\beta = -0.935$).

\item \textbf{Mechanism: a specific innate-immune axis, with a pharmacological rescue.}
Yan et al.\ close the loop from 21-day PM$_{2.5}$ exposure through mitochondrial damage
and mtDNA release to cGAS--STING activation and macrophage M1 polarisation, and show
that the cGAS inhibitor RU.521 attenuates the whole chain in vivo.

\item \textbf{One estimate, not zero --- but it is small and it is a mixture.} After
yesterday's zero-effect issue, the forest plot is populated again, with a single row.
See the note on the figure: the axis is now floored at a 1.6$\times$ span precisely so a
1.1\% risk ratio is not drawn as though it were a large one.
\end{itemize}

\clearpage

% ==================================================================== ANALYTICS
\section*{\sffamily\bfseries\color{Deep}Corpus analytics}
\vspace{-1.5mm}{\color{Amber}\rule{\linewidth}{1.1pt}}\vspace{3mm}

\noindent\includegraphics[width=\linewidth]{\FIGDIR/f1_subtopics.png}
\figcap{Subtopic assignment of the 19 in-scope records; single-label by dominant
contribution. \textit{Exposure assessment \& modelling} takes 8 of 19, more than twice
any other cluster, and \textit{Burden, policy \& mitigation} takes 4. Three of the ten
canonical clusters are empty today --- cardiovascular/metabolic, respiratory/allergic and
reproductive/developmental --- which is the widest health-side gap since 8 August.}

\vspace{2mm}
\noindent\includegraphics[width=\linewidth]{\FIGDIR/f2_design_endpoint.png}
\figcap{Left --- study architecture. Measurement campaigns (6) and modelling/inventory
work (5) together account for 11 of 19; primary human epidemiology contributes only 3
(two cohorts, one ecological analysis). Two records are \textit{Metadata only}: they were
surfaced through the Crossref journal sweep with no abstract deposited, so design,
sample and endpoint could not be extracted and are not guessed. Right --- health endpoint
by organ system. Eleven of nineteen records carry no human health endpoint at all, which
follows directly from the exposure/deposition weighting on the left.}

\vspace{2mm}
\noindent\includegraphics[width=\linewidth]{\FIGDIR/f3_metric_geography.png}
\figcap{Left --- particle metric. PM$_{2.5}$ alone is the exposure in 8 records; another 8
are unspeciated PM or emissions, which here mostly means deposited or wind-eroded
material rather than a regulated size fraction. Only one record is composition-resolved
in the register sense, though three more (Wei, Chen, Chai) are component-resolved in
their analysis. Right --- geographic provenance. Europe leads with 6 and China follows
with 5; sub-Saharan Africa contributes nothing today, and the two Latin American records
are both Mexican or Brazilian modelling rather than measurement.}

\vspace{2mm}
\noindent\includegraphics[width=\linewidth]{\FIGDIR/f6_lifecourse.png}
\figcap{Life-course windows addressed, counted only where a record's study population is
explicitly anchored at that stage. Two of the five windows are occupied and three are
empty: the day produced no paediatric, adolescent or in-utero record at all. The two
that are occupied sit at opposite ends --- a working-age UK Biobank cohort and a
$\geq$65 SEER-Medicare cohort.}

\vspace{2mm}
\noindent\includegraphics[width=\linewidth]{\FIGDIR/f4_heatmap.png}
\figcap{Subtopic $\times$ study-architecture cross-tabulation. The \textit{Exposure
assessment} row spreads across four architectures, which is what a broad but shallow
day looks like. The two cohort cells (Neuro, Other clinical) are the only cells in the
table that carry a human health endpoint together with a quantitative exposure contrast.}

\vspace{2mm}
\noindent\includegraphics[width=\linewidth]{\FIGDIR/f5_forest.png}
\figcap{The single forest-plottable estimate in the corpus. \textbf{The axis span is
enforced at a minimum of 1.6$\times$ in log space, and this issue is the reason.} With
one near-null estimate the axis previously auto-scaled to the data --- a range of 1.5\%
--- and drew RR 1.011 two-thirds of the way from the null line to the right edge. Nothing
was false and everything read as a large effect. Chen et al.'s PM$_{2.5}$--depression
result is deliberately \emph{not} plotted: the abstract reports ``HRs up to 1.20'' across
six exposure variables with no interval attached, and placing an upper envelope on a
shared risk axis would require inventing a confidence interval.}

\clearpage

% ==================================================================== DIGEST
\section*{\sffamily\bfseries\color{Deep}Paper-level digest}
\vspace{-1.5mm}{\color{Amber}\rule{\linewidth}{1.1pt}}\vspace{1mm}

{\fontsize{7.4}{9}\selectfont\color{Slate}Ordered by cluster. Each entry gives the
methodological core and the finding that matters, not an abstract paraphrase.\par}

\band{Exposure assessment \& modelling\hfill 8 records}{Sky}

\paperentry{Chai et al.\ 2026}
{Role of ammonia, aerosol pH and liquid water content in secondary inorganic aerosol formation in suburban Nanjing}
{Environ Pollut}
{Hourly water-soluble ions and precursor gases, January--March 2021, coupled to
ISORROPIA-II and WRF-Chem. SIA is \textbf{49.55\% of PM$_{2.5}$ mass} (mean
34.4\,$\mu$g/m$^3$), nitrate-dominated. Continuous ammonia excess fully neutralises the
anion pool; rising aerosol liquid water enhances aqueous oxidation and drives nitrate and
ammonium into the particle phase, while pH governs both semi-volatile partitioning and
the sulfate pathway (transition-metal-ion catalysis dominant). Thermodynamically the
system is \textbf{total-nitrate-limited}: PM$_{2.5}$ responds minimally to moderate
NH$_3$ cuts and needs $>$60\% reduction to move, whereas cutting HNO$_3$ alone matches
joint HNO$_3$+NH$_3$ control. The policy inversion --- prioritise NO$_x$, treat NH$_3$ as
a local adjunct --- is the transferable result.}
{10.1016/j.envpol.2026.128998}

\paperentry{Kan et al.\ 2026}
{Short-term observations of airborne microplastics in Phnom Penh: concentrations, aerodynamic size distribution and polymer composition}
{Atmosphere}
{First AMP observations for Cambodia, by $\mu$FTIR-ATR imaging, with meteorology,
gaseous co-pollutants, water-soluble ionic tracers and HYSPLIT back-trajectories.
Four independent 72-h periods: \textbf{0.97 $\pm$ 0.30 MP\,m$^{-3}$ in TSP} (range
0.55--1.27) and \textbf{0.33 $\pm$ 0.10 MP\,m$^{-3}$ in PM$_{2.5}$} (0.23--0.49);
134 particles identified in TSP, 46 of them inside the PM$_{2.5}$ fraction. PE, PP and
PET dominate, 52\% fragments, 82\% below 30\,$\mu$m Feret diameter. The authors are
explicit that $n = 4$ precludes source attribution --- read the polymer split as a
baseline, not an apportionment.}
{10.3390/atmos17080804}

\paperentry{Joshi et al.\ 2026 (preprint)}
{Environmental drivers of foliar particulate accumulation, leaf structure and biochemistry in urban trees in Delhi}
{Research Square}
{Ten native species, four contrasting Delhi sites, summer and winter 2025, with
branch-level triplicates. Total leaf-deposited PM rose from \textbf{120.13 $\pm$ 30.99 to
423.69 $\pm$ 30.99\,$\mu$g\,cm$^{-2}$} between seasons, fine PM from 6.82 $\pm$ 1.32 to
25.75 $\pm$ 1.32. \textit{Plumeria alba} retained most (414.92), \textit{Azadirachta
indica} least (140.82). Higher loads tracked ascorbic acid and foliar pH positively and
carotenoids negatively, and were \emph{not} explained by specific leaf area, leaf dry
matter content or thickness; APTI showed no seasonal change. The useful negative result
is that capture capacity and physiological tolerance are separate axes, so single-metric
species screening is inadequate.}
{10.21203/rs.3.rs-10524655/v1}

\paperentry{Alaimo et al.\ 2026}
{Elemental composition and Pb isotopic signatures in pine needles from the industrialized Milazzo area}
{Atmosphere}
{Forty \textit{Pinus pinea} needle samples by ICP-OES/ICP-MS across 36 elements, 25 of
them also for $^{206}$Pb/$^{207}$Pb and $^{208}$Pb/$^{206}$Pb. Enrichment factors above
10 for Cd, Cu, Mo, Sb, V and Zn against EF 0.5--2 for Al, Fe and Ti cleanly separate
anthropogenic from geogenic contributions. Pb ratios of 1.153--1.192 and 2.063--2.108
indicate mixing of industrial emissions with local geological background and, notably,
\emph{limited} residual leaded-gasoline signal. A good template for isotope-assisted
source apportionment where no monitoring network exists.}
{10.3390/atmos17080805}

\paperentry{Lazo et al.\ 2026}
{Biomonitoring trace metal deposition with \textit{Ligustrum lucidum}: Albanian fire incidents}
{J Environ Qual}
{Passive biomonitoring across 14 urban, industrial and rural sites in Elbasan, leaves
collected two weeks after an industrial-waste fire and a landfill fire. Al, V, Fe, Cu,
Cd, Pb and Zn showed low-to-moderate variability; Cr, Co, Ni, As and Sb were highly
variable with industrial-site outliers, and cluster analysis separated the two groups.
Method demonstration rather than exposure assessment --- there is no control period and
no airborne concentration --- but it is the kind of post-incident evidence that is
otherwise unobtainable.}
{10.1002/jeq2.70243}

\paperentry{Yang et al.\ 2026}
{Impacts of anthropogenic aerosols on a snowfall event in the Guanzhong Basin and surrounding areas}
{Atmos Chem Phys}
{Cloud-resolving fully coupled WRF-Chem, with emission-perturbation experiments
separating aerosol--radiation from aerosol--cloud interactions. The precipitation
response reverses sign between the basin and its surroundings: inside the basin,
excluding ARI gives a slight precipitation increase with emissions via enhanced
ice-phase processes, and ARI itself increases precipitation once the emission increase
passes a threshold (higher relative humidity, larger ice water path, better survival of
falling ice); outside, precipitation falls with emissions as ACI suppresses the liquid
phase. Spatial distribution is largely insensitive. Relevant to anyone using
precipitation as a wet-deposition assumption in a polluted basin.}
{10.5194/acp-26-11923-2026}

\paperentry{Gon\c{c}alves et al.\ 2026}
{Extreme Saharan dust outbreak: air quality crisis in NW Spain (March 2022)}
{Atmos Pollut Res}
{\textbf{Metadata only.} Surfaced through the Crossref journal sweep with no abstract
deposited. Title, journal, DOI and authorship are verified; design, sample size and
quantitative outcome are not extracted and are not inferred. Flagged for abstract
retry on the next run.}
{10.1016/j.apr.2026.103183}

\paperentry{Zheng et al.\ 2026}
{Aerosol modulation of precipitation systems migrating from mountains to plains in Northern China during early autumn}
{Atmos Res}
{\textbf{Metadata only.} Same provenance and same caveat as the entry above. The topic
is adjacent to Yang et al.\ on this page, but nothing in the two can be compared until
an abstract or full text is available.}
{10.1016/j.atmosres.2026.109289}

\band{Burden, policy \& mitigation\hfill 4 records}{Coral}

\paperentry{Liao et al.\ 2026}
{Beyond mortality: multi-outcome assessment of health economic benefits from air pollution control in China}
{Environ Res}
{Daily mortality, hospital admissions and outpatient visits across 18 Guangdong cities,
2015--2022, with generalized additive concentration--response models and valuation by
value of statistical life plus direct medical cost. PM$_{2.5}$ and NO$_2$ control
delivered \textbf{10{,}305.35 and 1{,}017.00 million CNY} respectively; O$_3$ control
returned \emph{negative} net benefits on every outcome. The headline methodological
number is that multi-outcome benefits of PM$_{2.5}$ control \textbf{exceed mortality-only
estimates by 31.79\%} --- i.e.\ the standard mortality-only accounting understates the
case by roughly a third. Over 85\% of mortality-related benefit accrues to the elderly.}
{10.1016/j.envres.2026.125501}

\paperentry{Bosso et al.\ 2026}
{AIRTREE v2.0: predicting ecosystem services provided by Italian urban parks under climate change scenarios}
{J Environ Manage}
{Multi-layer canopy model over 43 national-reforestation-programme sites under RCP4.5
and RCP8.5 to 2054, calibrated against five ICOS flux sites (Kling-Gupta efficiency up
to 0.90). Under RCP8.5 net primary productivity rises 9\% and evapotranspiration 8\%,
O$_3$ and NO$_2$ removal fall, and \textbf{PM$_{10}$ removal increases slightly}.
Species responses diverge sharply --- \textit{Pinus halepensis} leads on PM removal,
\textit{Sorbus torminalis} on O$_3$ uptake, \textit{Quercus pubescens} on carbon.
Open-access tool; the species-resolved PM output is the part worth reusing.}
{10.1016/j.jenvman.2026.130738}

\paperentry{Fares et al.\ 2026}
{Can urban green infrastructure influence the urban heat island effect and air quality nexus?}
{Environ Pollut}
{Review arguing that the thermal and air-quality benefits of green infrastructure are
almost never assessed together, and synthesising the interaction: shading and
evapotranspiration lower leaf surface temperature, which promotes thermophoretic
deposition of particles, alongside foliar uptake of gaseous pollutants and inhibition of
secondary aerosol formation. Well-maintained greenspace outperforms engineered
installations, whose effects are more localised. The limitations section is the honest
part --- BVOC emission, small effect sizes, species and climate-zone dependence, and
canopy-induced night-time warming.}
{10.1016/j.envpol.2026.128960}

\paperentry{Moghaddam et al.\ 2026 (preprint)}
{Reducing particulate matter emissions in landfills using novel dust suppressants}
{Research Square}
{Wind-tunnel factorial on soil from the Meyami landfill in Mashhad: three mesh sizes,
four suppressants (compost, vermicompost, bentonite, nano-clay bentonite), three doses,
two exposure times, two wind speeds. Mesh 20 was optimal across all suppressants
($>$85\% reduction in wind-erosion rate); \textbf{compost at 10\% dose reached 86.52\%}
and bentonite at 20\% reached 80\%, both at 4\,m/s. A bench result on a single soil ---
no field validation and no airborne concentration measured --- but the compost finding is
operationally cheap enough to be worth a field trial.}
{10.21203/rs.3.rs-10431018/v1}

\clearpage

\band{Sensing, forecasting \& instrumentation\hfill 2 records}{Amber}

\paperentry{Ng et al.\ 2026}
{Impacts of droughts on biomass burning emissions, air quality and public health in the Amazon}
{Atmos Chem Phys}
{Satellite fire activity plus GFEDv4s emissions driving GEOS-Chem High Performance,
2010--2015, with paired fire-on/fire-off simulations. In non-drought years biomass
burning contributes \textbf{34.0\% (32.0--35.5) of dry-season CO and PM$_{2.5}$}; in
drought years that rises to \textbf{55.4\% (44.1--76.7)}. Fire-induced PM$_{2.5}$ and
ozone raise premature mortality (GEMM) by 6.0\% and 18.6\% in non-drought years, 8.9\%
and 24.4\% in drought years --- \emph{at stable deforestation rates}, which is the point:
the climate-driven amplification is separable from the land-use driver. For network
design, the drought-year contribution range is wide enough that a fixed seasonal
baseline will misestimate the fire fraction by a factor of two.}
{10.5194/acp-26-11835-2026}

\paperentry{Legaria-Santiago et al.\ 2026}
{Leveraging remote traffic data for local air pollutant estimation across London monitoring sites}
{Atmosphere}
{Four interpretable tree ensembles (Random Forest, Extra Trees, LightGBM, XGBoost)
across six predictor scenarios, from remote traffic + meteorology + temporal variables
alone up to including one and four neighbouring stations, benchmarked against ridge
regression, spatial interpolation and a cross-site validation. Without neighbouring
stations, NO$_2$ RMSE fell from \textbf{9.73--11.66 to 8.72--11.52\,$\mu$g/m$^3$} when
traffic information was added, and SHAP shows traffic variables contributing at a level
comparable to neighbouring-station measurements in traffic-dominated settings. The
practical claim --- remote traffic counts can partially substitute for network density ---
is demonstrated for NO$_2$; the PM channels are reported but not where the gain is.}
{10.3390/atmos17080806}

\band{Other clinical endpoints\hfill 2 records}{Slate}

\paperentry{Wei et al.\ 2026}
{Chemical composition of fine particulate matter and all-cause mortality of older lung cancer patients}
{Environ Res}
{528{,}127 SEER-Medicare adults aged $\geq$65 diagnosed with lung cancer 2000--2019,
followed annually; two-year moving-average exposure to 15 PM$_{2.5}$ components from
high-resolution spatiotemporal models linked by residential ZIP code; generalized
weighted quantile sum regression with random holdouts, adjusted for demographics,
histology, stage, first-course treatment, comorbidity and neighbourhood covariates.
\textbf{RR 1.011 (95\% CI 1.010--1.013) per decile increase across all components.}
Contribution differences were modest but silicon, nitrate, vanadium, zinc and iron led
--- road dust, traffic, fossil-fuel and heavy-fuel-oil combustion. Exploratory subgroups
suggest a stronger association in non-small-cell, later-stage and female patients.}
{10.1016/j.envres.2026.125538}

\paperentry{Guerrero-Morales et al.\ 2026}
{Spatial distribution of long-term ultrafine particle concentrations and systemic lupus erythematosus burden in central Mexico}
{Environ Toxicol Pharmacol}
{Municipal-level ecological analysis: SLE mortality 1998--2023 and hospitalisations
2008--2024 from national databases against satellite-model UFP for 2010--2019, using rate
ratios, Jenks natural breaks and principal component analysis. National SLE mortality
0.23--0.86 per 100{,}000; municipal mean UFP roughly
\textbf{1{,}500--24{,}000 particles/cm$^3$}. Higher UFP and higher hospitalisation
burden co-occur in metropolitan municipalities; mortality is heterogeneous and the
authors say it should be read descriptively. Ecological design, no individual exposure,
explicitly hypothesis-generating --- but it is one of very few UFP--autoimmunity datasets
from a middle-income country.}
{10.1016/j.etap.2026.105143}

\band{Neuro / mental health\hfill 1 record}{Deep}

\paperentry{Chen et al.\ 2026}
{Integrated plasma proteomics and metabolomes reveal potential biomolecules linking PM$_{2.5}$ to depression}
{Environ Toxicol Pharmacol}
{24{,}940 depression-free UK Biobank participants, exposure to PM$_{2.5}$ and five
constituents (ammonium, nitrate, sulfate, elemental carbon, organic carbon), Olink
proteomics and NMR metabolomics, Cox models with mediation, functional enrichment and
drug-repurposing prediction. \textbf{982 incident cases over a mean 13 years}; HRs
reported ``up to 1.20 per SD'' with ammonium the strongest constituent. \textbf{140
proteins and 3 metabolites} were statistically significant mediators, mediation
proportions up to 8.00\% and 3.42\%, enriched in lipid and unsaturated fatty acid
metabolism, immune regulation and inflammation. Note that the abstract gives no interval
for the 1.20 figure, which is why it does not appear on this issue's forest plot.}
{10.1016/j.etap.2026.105144}

\band{Mechanistic toxicology\hfill 1 record}{Violet}

\paperentry{Yan et al.\ 2026}
{PM$_{2.5}$ induces lung injury via mtDNA-cGAS-STING-mediated macrophage M1 polarization}
{Front Pharmacol}
{21-day subacute PM$_{2.5}$ mouse model plus bone-marrow-derived macrophages in vitro.
Exposure produced pulmonary inflammation and tissue injury with macrophage accumulation
and M1 polarisation; mechanistically, mitochondrial damage released mtDNA and activated
cGAS--STING. The load-bearing experiment is the rescue: \textbf{pharmacological cGAS
inhibition with RU.521 attenuated STING activation, M1 polarisation and lung injury in
vivo}, which converts a correlative pathway into an interventional one. Subacute rather
than chronic, and a single exposure route --- the dosimetry is the usual weak point.}
{10.3389/fphar.2026.1902928}

\band{Occupational \& indoor\hfill 1 record}{Clay}

\paperentry{Liu \& Hu 2026 (preprint)}
{Multivariable determinants of indoor PM$_{2.5}$ and infiltration dynamics in 48 residential buildings across eight cities of Northern China}
{Preprints.org}
{Paired indoor/outdoor PM$_{2.5}$ at 5-minute resolution for seven consecutive days in
each of four seasons across Beijing, Tianjin, Shijiazhuang, Taiyuan, Jinan, Zhengzhou,
Xi'an and Anyang --- \textbf{32{,}256 hourly observations from 48 residences}, with
building characteristics, occupant behaviour and meteorology recorded alongside. A
multivariable OLS model with cluster-robust standard errors ($R^2 = 0.879$) gives
outdoor PM$_{2.5}$ $\beta = 0.891$, window-open fraction $\beta = 0.572$, cooking
$\beta = 0.049$ and \textbf{air-purifier operation $\beta = -0.935$}, all $p < 0.001$.
Infiltration factor \textbf{0.28 $\pm$ 0.15}; outdoor sources contribute
\textbf{57--69\%} of indoor PM$_{2.5}$, peaking in the heating season; median deposition
rate 0.53\,h$^{-1}$ against a median air change rate of 0.52\,h$^{-1}$, exceeding
2\,h$^{-1}$ with windows open. Preprint, not peer reviewed --- but the sample size and
the seasonal design are both well beyond what this literature usually carries.}
{10.20944/preprints202608.1362.v1}

\vspace{4mm}
% ==================================================================== TABLE
\section*{\sffamily\bfseries\color{Deep}Corpus register}
\vspace{-1.5mm}{\color{Amber}\rule{\linewidth}{1.1pt}}\vspace{2.5mm}

{\fontsize{6.7}{8.4}\selectfont
\setlength{\tabcolsep}{2.4pt}
\begin{longtable}{@{}>{\RaggedRight}p{25mm}>{\RaggedRight}p{41mm}>{\RaggedRight}p{27mm}%
>{\RaggedRight}p{23mm}>{\RaggedRight}p{18mm}>{\RaggedRight}p{15mm}@{}}
\toprule
\textsf{\textbf{First author}} & \textsf{\textbf{Focus}} & \textsf{\textbf{Journal}} &
\textsf{\textbf{Design}} & \textsf{\textbf{Metric}} & \textsf{\textbf{Region}}\\
\midrule
\endfirsthead
\toprule
\textsf{\textbf{First author}} & \textsf{\textbf{Focus}} & \textsf{\textbf{Journal}} &
\textsf{\textbf{Design}} & \textsf{\textbf{Metric}} & \textsf{\textbf{Region}}\\
\midrule
\endhead
\multicolumn{6}{@{}l@{}}{\cellcolor{Mist}\textsf{\textbf{\textcolor{Deep}{Neuro / mental health}}}}\\[0.6mm]
\href{https://doi.org/10.1001/jamanetworkopen.2026.25015}{Dales} & Wildfire-sourced PM$_{2.5}$ and ED visits for migraine and other primary headache & \textit{JAMA Netw Open} & Case-crossover & PM$_{2.5}$ (wildfire) & Canada\\
\href{https://doi.org/10.1093/ageing/afag215}{Glover} & Umbrella meta-review of urban environmental exposures and cognitive health & \textit{Age Ageing} & Umbrella review / meta-analysis & PM$_{2.5}$ & Global\\
\href{https://doi.org/10.1016/j.nsa.2026.107019}{Goossens} & In utero air pollution exposure and pre-/postnatal brain development & \textit{Neurosci Appl} & Narrative review & PM$_{2.5}$ / PM$_{10}$ / PM$_{\mathrm{coarse}}$ & Global\\
\href{https://doi.org/10.1016/j.isci.2026.116617}{Ma} & PM$_{2.5}$-induced neurotoxicity and depression: exposure risks and mechanisms & \textit{iScience} & Narrative review & PM$_{2.5}$ + components & Global\\
\href{https://doi.org/10.3760/cma.j.cn112150-20250804-00752}{Pan} & Maternal PM exposure and offspring neurodevelopment: effects and mechanisms & \textit{Zhonghua Yu Fang Yi Xue Za Zhi} & Narrative review & PM (bulk) & China\\
\addlinespace[0.9mm]
\multicolumn{6}{@{}l@{}}{\cellcolor{Mist}\textsf{\textbf{\textcolor{Teal}{Cardiovascular \& metabolic}}}}\\[0.6mm]
\href{https://doi.org/10.3760/cma.j.cn112338-20251120-00843}{Huang} & PM$_{2.5}$ and dyslipidaemia in rural Chinese children; dietary effect modification & \textit{Zhonghua Liu Xing Bing Xue Za Zhi} & Cross-sectional (multistage) & PM$_{2.5}$ & China\\
\href{https://doi.org/10.32394/pe/218966}{Kazmierska} & Air pollution and cardiovascular disease: review of current data & \textit{Przegl Epidemiol} & Narrative review & PM$_{2.5}$ / PM$_{10}$ & Global\\
\href{https://doi.org/10.1007/s00125-026-06814-2}{Xi} & Air pollution, built environment and temperature vs. prediabetes phenotypes & \textit{Diabetologia} & Prospective cohort & PM$_{2.5}$ & Germany\\
\href{https://doi.org/10.1161/JAHA.125.046917}{Yang} & Short-term PM$_{2.5}$ and hospitalisation/ED use in atrial fibrillation & \textit{J Am Heart Assoc} & Case-crossover & PM$_{2.5}$ (wildfire-influenced) & USA\\
\addlinespace[0.9mm]
\multicolumn{6}{@{}l@{}}{\cellcolor{Mist}\textsf{\textbf{\textcolor{Sky}{Reproductive \& developmental}}}}\\[0.6mm]
\href{https://doi.org/10.1016/j.jes.2025.10.043}{Chen} & Prenatal air pollution, fetal growth trajectory phenotypes and child neurodevelopment & \textit{J Environ Sci (China)} & Prospective cohort & PM$_{2.5}$ / PM$_{10}$ & China\\
\href{https://doi.org/10.1016/j.ecoenv.2026.120566}{Huang} & Joint prenatal bisphenol + air-pollutant mixture exposure and preschool IQ & \textit{Ecotoxicol Environ Saf} & Prospective cohort (mixtures) & PM$_{2.5}$ / PM$_{10}$ & China\\
\href{https://doi.org/10.1007/s40572-026-00552-8}{Peterson} & Ambient air pollution, PAHs and pubertal development: critical review & \textit{Curr Environ Health Rep} & Critical review & PM$_{2.5}$ components & Global\\
\href{https://doi.org/10.1016/j.envres.2026.125313}{Xie} & Long-term PM$_{2.5}$ and timing of pubertal development in boys and girls & \textit{Environ Res} & Prospective cohort & PM$_{2.5}$ & China\\
\addlinespace[0.9mm]
\multicolumn{6}{@{}l@{}}{\cellcolor{Mist}\textsf{\textbf{\textcolor{Sage}{Respiratory \& allergic}}}}\\[0.6mm]
\href{https://doi.org/10.1007/s00484-026-03289-z}{Bayar Muluk} & Climate change and otolaryngology: airway disease impacts and emission reduction & \textit{Int J Biometeorol} & Narrative review & PM$_{2.5}$ & Global\\
\href{https://doi.org/10.4103/mgr.MEDGASRES-D-25-00164}{Liao} & Short-term air pollution, respiratory symptoms and lung function in asthma & \textit{Med Gas Res} & Panel study & PM$_{2.5}$ / PM$_{10}$ & China\\
\href{https://doi.org/10.3390/epidemiologia7040097}{Sanchez de Prada} & Ambient PM and asthma-coded admissions in children 0-3 y, Spain & \textit{Epidemiologia} & Ecological panel & PM$_{2.5}$ / PM$_{10}$ & Spain\\
\href{https://doi.org/10.1016/j.jhazmat.2026.143076}{Yang} & Allergenic vs. toxic risk of airborne seed hairs compared with ambient PM$_{2.5}$ & \textit{J Hazard Mater} & Field measurement & PM$_{2.5}$ + bioaerosol & China\\
\addlinespace[0.9mm]
\multicolumn{6}{@{}l@{}}{\cellcolor{Mist}\textsf{\textbf{\textcolor{Violet}{Mechanistic toxicology}}}}\\[0.6mm]
\href{https://doi.org/10.1016/j.jes.2025.12.074}{Ding} & PM$_{2.5}$ promotes vascular calcification via miR-27b-5p/HES1 signalling & \textit{J Environ Sci (China)} & In vitro + organoid + murine & PM$_{2.5}$ & China\\
\href{https://doi.org/10.1007/s00360-026-01698-5}{Maraschi} & Osmoregulatory disruption and metal bioaccumulation from settleable atmospheric PM & \textit{J Comp Physiol B} & Controlled exposure (ecotox) & Settleable APM & Brazil\\
\href{https://doi.org/10.1016/j.jes.2025.11.018}{Wu} & Macrophage NLRP3 inflammasome in real-ambient PM-induced hepatic glucose dysregulation & \textit{J Environ Sci (China)} & Real-ambient murine + co-culture & Real-ambient PM & China\\
\addlinespace[0.9mm]
\multicolumn{6}{@{}l@{}}{\cellcolor{Mist}\textsf{\textbf{\textcolor{Clay}{Occupational \& indoor}}}}\\[0.6mm]
\href{https://doi.org/10.3389/fmedt.2026.1796828}{Chen} & Nested suction instrument to cut intraoperative surgical-smoke PM$_{2.5}$ & \textit{Front Med Technol} & Bench / device experiment & PM$_{2.5}$ (surgical smoke) & China\\
\href{https://doi.org/10.1093/annweh/wxag055}{Cole} & Determinants of respirable crystalline silica exposure in Australian tunnelling & \textit{Ann Work Expo Health} & Observational exposure study & Respirable dust / RCS & Australia\\
\href{https://doi.org/10.1016/j.jhazmat.2026.143081}{Ding} & Stage-resolved hair Zn-isotope dynamics as a firefighting combustion-exposure tracer & \textit{J Hazard Mater} & Repeated-measures biomarker & Combustion-derived PM & China\\
\href{https://doi.org/10.3390/jox16040119}{Ramsden} & Xenobiotic hazards in aircraft cabin air & \textit{J Xenobiot} & Review + risk economics & Ultrafine particles & Global\\
\addlinespace[0.9mm]
\multicolumn{6}{@{}l@{}}{\cellcolor{Mist}\textsf{\textbf{\textcolor{Amber}{Sensing, forecasting \& instrumentation}}}}\\[0.6mm]
\href{https://doi.org/10.1016/j.jenvman.2026.130515}{Shi} & Adaptive spatiotemporal graph transformer for multi-site multi-step PM$_{2.5}$ forecasting & \textit{J Environ Manage} & Deep-learning model & PM$_{2.5}$ & China / India\\
\href{https://doi.org/10.3390/bios16070380}{Tsuchiya} & Seasonal and indoor-outdoor characteristics of fluorescent aerosol particles in offices & \textit{Biosensors} & Real-time sensor campaign & Bioaerosol / coarse & Japan\\
\addlinespace[0.9mm]
\multicolumn{6}{@{}l@{}}{\cellcolor{Mist}\textsf{\textbf{\textcolor{Amber}{Exposure assessment \& modelling}}}}\\[0.6mm]
\href{https://doi.org/10.1029/2025GH001621}{Solorzano-Araque} & Multiscale exposure modelling: LOTOS-EUROS + mobility surveys + particle toxicity & \textit{GeoHealth} & CTM + mobility integration & PM$_{2.5}$ / PM$_{10}$ & Colombia\\
\addlinespace[0.9mm]
\multicolumn{6}{@{}l@{}}{\cellcolor{Mist}\textsf{\textbf{\textcolor{Coral}{Burden, policy \& mitigation}}}}\\[0.6mm]
\href{https://doi.org/10.1002/hsr2.72886}{Karamzad} & Ischaemic heart disease burden in MENA 1990-2021; PM attributable fraction 28.2\% & \textit{Health Sci Rep} & GBD secondary analysis & Ambient PM & MENA\\
\href{https://doi.org/10.1371/journal.pone.0347989}{Rees} & HEPA filtration to reduce respiratory infection in care homes & \textit{PLoS One} & RCT process evaluation & Filtered PM & UK\\
\href{https://doi.org/10.1016/j.jes.2026.01.009}{Shan} & Policy-driven emission reductions and post-pandemic rebound, high-resolution inventories & \textit{J Environ Sci (China)} & Emission inventory & PM emissions & China\\
\addlinespace[0.9mm]
\multicolumn{6}{@{}l@{}}{\cellcolor{Mist}\textsf{\textbf{\textcolor{Slate}{Other clinical endpoints}}}}\\[0.6mm]
\href{https://doi.org/10.3390/dj14070440}{Lupita} & PM$_{2.5}$, preventive-care capacity and caries/edentulism burden across the EU & \textit{Dent J} & Ecological panel & PM$_{2.5}$ & EU-27\\
\href{https://doi.org/10.1016/j.jdin.2026.06.008}{Park} & Long-term PM exposure and risk of androgenetic alopecia & \textit{JAAD Int} & Nationwide cohort & PM$_{2.5}$ / PM$_{10}$ & South Korea\\
\href{https://doi.org/10.3171/2026.3.PEDS2615}{Sun} & Environmental exposure and area deprivation vs. paediatric hydrocephalus outcomes & \textit{J Neurosurg Pediatr} & Registry cohort + geospatial & PM$_{2.5}$ (satellite-derived) & USA\\
\addlinespace[0.9mm]
\bottomrule
\end{longtable}}


\vspace{4mm}
\begin{keybox}
{\sffamily\bfseries\footnotesize\color{Deep}Provenance and method}\\[1.2mm]
{\fontsize{7.2}{8.8}\selectfont
Entry window \textbf{2026-08-21 $\rightarrow$ 2026-08-21}, contiguous with the 20 August
issue. Built at 22:1x local on 22 August under the 22:00 rule, as a one-day backfill.
\textbf{Sources.} PubMed E-utilities on \texttt{[EDAT]} returned 10 unique records
(6 shipped). Europe PMC on \texttt{CREATION\_DATE} returned 28 unique records, of which
\textbf{14 were archival PMC back-scans} --- legacy scans with no title and a publication
year decades outside the window, re-ingested with a creation date inside it --- 7 were
already in \texttt{seen.json}, and 6 shipped. The Crossref by-ISSN journal sweep returned
34 records and shipped 7. \textbf{OpenAlex returned HTTP 429 for the 19th consecutive
run} and remains dead. arXiv returned nothing new under the 3-day recency screen.
\textbf{Screening.} 46 records were screened out and written to
\texttt{state/rejected.jsonl} with an individual reason each; the largest single reason
is the Europe PMC back-scan class, followed by the Crossref sweep's non-airborne
chemistry (PFAS, microplastics in water, sediment geochemistry) and four retraction
notices. Two borderline exclusions are logged as borderline: cross-floor bioaerosol
transmission in drainage stacks, and an occupant-comfort building-physics cluster.
\textbf{Labels.} Subtopic, design, particle-metric and geography are assigned from title,
abstract, keywords and MeSH; assignment is single-label by dominant contribution and is a
judgement, not a controlled vocabulary. Two records carry no abstract at all and are
marked \textit{Metadata only} rather than being characterised by inference.
Effect estimates are transcribed verbatim from abstracts.
All DOIs were checked by \texttt{check\_dois.py} before publication. Content derived from
publisher metadata; abstracts remain the copyright of their respective publishers.\par}
\end{keybox}

\end{document}
